% Challenge dossier: VI.06.C3
\subsection*{Challenge VI.06.C3 --- When does a prime stay prime after extending the field?}
\label{chal:vi-06-c3}

\paragraph{Problem.}
Let \(k\subseteq K\) be a field extension and let \(f\in k[t]\) be irreducible.  The ideal
\((f)\subset k[t]\) is prime.  Investigate the extended ideal \((f)\subset K[t]\):
\begin{enumerate}
\item prove that it is prime exactly when \(f\) remains irreducible in \(K[t]\);
\item describe what happens when \(f\) factors into distinct irreducibles over \(K\);
\item analyze \(f=t^2+1\) for \(\mathbb R\subset\mathbb C\).
\end{enumerate}

\ifdefined\FullProblemDossiers
\paragraph{Why this problem matters.}
A prime geometric point over one field may split after base change.  This is the elementary
one-variable model for geometric irreducibility.

\paragraph{Useful ideas before starting.}
The quotient-domain criterion, factorization in a PID, and the Chinese remainder theorem
for pairwise coprime factors.

\paragraph{Guided hints.}
Use \(K[t]/(f)\).  If \(f=f_1\cdots f_r\) with distinct irreducible factors, the quotient
has one minimal prime for each factor and, when the factors are pairwise coprime, decomposes
through CRT.

\paragraph{Solution.}
The extended ideal is prime exactly when \(K[t]/(f)\) is a domain, which in the PID
\(K[t]\) is equivalent to irreducibility of \(f\).  If
\[
f=f_1\cdots f_r
\]
with distinct irreducible factors, then the minimal primes over \((f)\) are \((f_i)\).  If
the factors are pairwise coprime, CRT gives
\[
K[t]/(f)\cong\prod_i K[t]/(f_i).
\]
For \(t^2+1\) over \(\mathbb C\), the factors are \(t-i\) and \(t+i\), so the single real
prime becomes two complex maximal primes.

\paragraph{What happened mathematically.}
Primality is relative to the coefficient field.  The future language of base change will
track exactly how points and irreducible components transform.

\paragraph{Extensions.}
Study inseparable factorization phenomena in positive characteristic and distinguish
reducibility from nonreducedness after base change.

\paragraph{Provenance.}
Legacy field-extension counterexample retained as an advanced preview; systematic treatment
is reserved for VI/24.
\fi
